\chapter{Ontbinden in factoren}
\label{chap:factoren}

\section{Ontbinden in factoren}
Het omgekeerde van haakjes wegwerken noemen we \emph{ontbinden in 
factoren}. Hierbij wordt een wiskundige uitdrukking waarin geen haakjes 
staan, omgezet in een \emph{product} van twee (of meer) factoren, en dus 
in een uitdrukking \emph{met} haakjes. Waar haakjes wegwerken het 
vereenvoudigen van een formule is, ``bouwen'' we bij ontbinden de haakjes 
juist \emph{terug}:
\[
\underbrace{3(x+4)}_{\text{product}} 
\quad 
\xrightleftharpoons[\text{ontbinden in factoren}]{\text{haakjes wegwerken}} 
\quad 
\underbrace{3x + 12}_{\text{som}}
\]

Ontbinden in factoren lijkt misschien een omwegen, maar het is een van de 
belangrijkste gereedschappen in de wiskunde. Zo kun je met ontbinden 
\emph{vergelijkingen} oplossen die je anders niet aankunt: als 
$x^2 - 5x = 0$ herschreven wordt als $x(x-5)=0$, zie je in \'e\'en oogopslag 
dat $x=0$ of $x=5$ de oplossingen zijn. Ook bij het vereenvoudigen van 
breuken met variabelen is ontbinden onmisbaar.

Handig om te weten: je kunt je antwoord altijd controleren door de haakjes 
in je uitkomst weer weg te werken, volgens de regels van Hoofdstuk \ref{chap:haakjes}. Kom je terug bij de oorspronkelijke 
uitdrukking, dan is het ontbinden gelukt.

Bij het ``ontbinden in factoren'' zijn de volgende regels daarbij behulpzaam:

\regel{Regel 1}{ab+ac=a(b+c)}

De termen $a \cdot b$ en $a \cdot c$ bevatten beiden dezelfde factor $a$. Door deze factor $a$ buiten haken te halen wordt de vorm $ab+ac$ ontbonden in factoren.
In feite is deze regel het omgekeerde van regel 1 uit hoofdstuk \ref{chap:haakjes}.

Voorbeelden:
\begin{align*}
5x + 5y  &= 5(x + y) \\
6a + 6b  &= \\
15p - 3q &= 3(5p - q) \\
16x - 4y &= \\
ax + 2ay - 5a    &= a(x + 2y - 5) \\
5bx - 10by - 15b &= \\
x^2 + 2x &= x(x + 2) \\
y^2 - 3y &=
\end{align*}

\regel{Regel 2}{a^2-b^2=(a+b)(a-b)}
Met deze regel kan het verschil van twee kwadraten worden ontbonden in factoren.
In feite is deze regel het omgekeerde van het merkwaardig product 1 uit hoofdstuk \ref{chap:haakjes}.

Voorbeelden:
\begin{align*}
x^2 - 25   &= x^2 - 5^2 = (x + 5)(x - 5) \\
y^2 - 36   &= \\
4a^2 - 49  &= (2a)^2 - 7^2 = (2a + 7)(2a - 7) \\
25b^2 - 16 &=
\end{align*}

\regel{Regel 3}
{
    x^2+\mathrm{b}x+\mathrm{c} = (x+\mathbf{p})(x+\mathbf{q}) \quad 
    \begin{cases}
        \mathbf{p}+\mathbf{q}      &= \mathrm{b} \\
        \mathbf{p}\cdot \mathbf{q} &= \mathrm{c}
    \end{cases}
}

Met deze regel kan de kwadratische vorm $x^2+b\cdot x+c$ worden ontbonden in factoren.
Hiervoor zoeken we 2 onbekende (gehele) getallen $p$ en $q$, waarvoor geldt:
de som van $p$ en $q$ is gelijk aan $b$, het product van $p$ en $q$ is gelijk aan $c$.
Voor de onbekende getallen $p$ en $q$ geldt dus:

De verantwoording van deze regel krijgen we als we de haakjes in regel 3 wegwerken: 
\begin{align*}
    (x+p)(x+q)&=x^2+qx+px+pq \\
              &=x^2+(p+q)x+pq
\end{align*}
Zodat: $x^2+bx+c=x^2+(p+q)x+pq$
Hieruit kunnen we concluderen, dat $p+q$ overeenkomt met $b$, en dat $p\cdot q$ overeenkomt met $c$.

Opmerking: De methodiek van regel 3 is niet altijd mogelijk (zie voorbeeld 5)!

\subsection*{Voorbeelden:}
\begin{enumerate}
    \itemsep1em
    \item Ontbind in factoren: $x^2+5x+6$ \\

    Voor het ontbinden van deze vorm in factoren zoeken we dus 2 gehele getallen $p$ en $q$ waarvoor geldt: de som $p+q$ is gelijk aan 5, het product $p\cdot q$ is gelijk aan 6. \\
    
    Bij het zoeken naar deze getallen $p$ en $q$ beginnen we altijd met het product, omdat er maar enkele combinaties van 2 getallen zijn die vermenigvuldigd 6 opleveren. Daarna bepaalt de som van deze getallen, welke combinatie de juiste is. De methode in onderstaande tabel kan daarbij helpen.
    
    \begin{table}[H]
        \centering
        \renewcommand{\arraystretch}{1.4}
        \begin{tabular}{|l|l|}
            \hline
            \rowcolor{vakgray}
            \textbf{Het product van $p$ en $q$ is 6} & \textbf{De som van $p$ en $q$ is $+5$} \\ \hline
            $p = 1$ en $q = 6$      & Som $= 7$, klopt niet \\ \hline
            $p = -1$ en $q = -6$    & Som $= -7$, klopt niet \\ \hline
            $p = 2$ en $q = 3$      & Som $= 5$, klopt \\ \hline
            $p = -2$ en $q = -3$    & Som $= -5$, klopt niet \\ \hline
        \end{tabular}
    \end{table}
    Dus: $x^2+5x+6=(x+2)(x+3)$ 

    \item Ontbind in factoren: $y^2+10y+24$ 

	\item Ontbind in factoren: $x^2-2x-8$

    \begin{table}[H]
        \centering
        \renewcommand{\arraystretch}{1.4}
        \begin{tabular}{|l|l|}
            \hline
            \rowcolor{vakgray}
            \textbf{Het product van $p$ en $q$ is $-8$} & \textbf{De som van $p$ en $q$ is $-2$} \\ hline
            $p = 1$ en $q = -8$     & Som $= -7$, klopt niet \\ \hline
            $p = -1$ en $q = 8$     & Som $= 7$, klopt niet \\ \hline
            $p = 2$ en $q = -4$     & Som $= -2$, klopt \\ \hline
            $p = -2$ en $q = 4$     & Som $= 2$, klopt niet \\ \hline
        \end{tabular}
    \end{table}

    \item Ontbind in factoren: $x^2+3x-18$ 

    \item Ontbind in factoren: $x^2-2x+10$

    De tabel wordt nu:

    \begin{table}[H]
        \centering
        \renewcommand{\arraystretch}{1.4}
        \begin{tabular}{|l|l|}
            \hline
            \rowcolor{vakgray}
            \textbf{Het product van $p$ en $q$ is $10$} & \textbf{De som van $p$ en $q$ is $-2$} \\ \hline
            $p = 1$ en $q = 10$     & Som $= 11$, klopt niet \\ \hline
            $p = -1$ en $q = -10$   & Som $= -11$, klopt niet \\ \hline
            $p = 2$ en $q = 5$      & Som $= 7$, klopt niet \\ \hline
            $p = -2$ en $q = -5$    & Som $= -7$, klopt niet \\ \hline
        \end{tabular}
    \end{table}

    Conclusie: omdat we geen gehele getallen p en q kunnen vinden, faalt deze methode.
\end{enumerate}

\newpage
\section{Huiswerkopdrachten}
Hieronder vind je een aantal opgaven waarbij je uitdrukkingen in factoren moet ontbinden. Werk elke opgave netjes uit en controleer eventueel je antwoord door de factoren weer te vermenigvuldigen.

\subsection*{Opgave: Ontbind in factoren.}
\renewcommand{\arraystretch}{1.6}
\begin{tabular}{m{15em}@{\hspace{4em}}l}
    \text{a.}\hspace{0.5em} $9x + 3y$           & \text{k.}\hspace{0.5em} $x^2 + 7x + 10$    \\
    \text{b.}\hspace{0.5em} $6a + 18$           & \text{l.}\hspace{0.5em} $x^2 + 6x + 8$     \\
    \text{c.}\hspace{0.5em} $12x - 9y$          & \text{m.}\hspace{0.5em} $x^2 - 2x - 63$    \\
    \text{d.}\hspace{0.5em} $4x + 2y - 12$      & \text{n.}\hspace{0.5em} $y^2 - 5y - 6$     \\
    \text{e.}\hspace{0.5em} $x^2 + 2x$          & \text{o.}\hspace{0.5em} $a^2 - 16a + 64$   \\
    \text{f.}\hspace{0.5em} $t^2 + 12t$         & \text{p.}\hspace{0.5em} $x^2 + 4x - 12$    \\
    \text{g.}\hspace{0.5em} $a^2 - 16$          & \text{q.}\hspace{0.5em} $x^2 + 10x + 25$   \\
    \text{h.}\hspace{0.5em} $p^2 - 64$          & \text{r.}\hspace{0.5em} $p^2 - 16p + 36$   \\
    \text{i.}\hspace{0.5em} $9x^2 - 25$         & \text{s.}\hspace{0.5em} $x^2 + 12x + 36$   \\
    \text{j.}\hspace{0.5em} $16x^2 - 4$         &                                            \\
\end{tabular}